\section{模型评价、适用范围与防灾启示}\label{sec:discussion}
\subsection{从区域异常到源区机制的跨尺度解释}
因素族排序与机制计算共同提供了一条可检验的解释路径。气象分析中，土壤水分的季节背景位置较高，而事件前168小时的变化方向为下降；这一组合提示应同时关注环境状态的历史积累和近期演变。源区水作用模型则以储水、排水、连通性和有效应力为状态，表达裂隙内的水力响应。两者分别对应区域土层和源区水路，联系它们的关键是空间代表性及流路支持，而非要求同名变量具有相同趋势。

热状态也具有明确的时标分工。在一维半无限、初始均温的线性热扩散条件下，$\alpha=1.5\times10^{-6}$\,m$^2$/s的阶跃边界示例中，72小时后2\,m和5\,m深处的温度响应与表面阶跃之比分别约为0.0233和$1.43\times10^{-8}$。这一深度依赖把解释重点分成浅部短时热扰动、多年初始热状态以及连通水路的传热效应。气象异常说明边界变化，热响应计算进一步回答该变化能够在何种深度与时间支持上影响强度。

因此，本文将TOP5作为环境背景输入的优先集合，将条件临界关系作为状态到稳定性的物理映射。结构几何与既有损伤决定共同基础，温度、水及耦合项通过配对设计表达各自贡献。该组织方式既保留了多源数据的统计信息，又将机制解释落实到应力、温度、压力与时间尺度。

\subsection{数值敏感性与新增过程的解释收益}
传播模型比较中，守恒、初态投影、时间步及边界分别代表不同的数值问题。本文同时报告过程场和这些响应量，比较新增物理过程所解释的现象。表\ref{tab:common-numerics}集中列出具有代表性的量化结果，作为读取图像和到达括号的尺度参照。

\begin{table}[htbp]
\centering\small
\caption{代表性数值对照。各行沿用其原实验配置，表内不同路线的量不作跨配置精度排名；目标括号来自实际存帧、相对整体释放。}\label{tab:common-numerics}
\begin{tabularx}{\linewidth}{p{2.7cm}Yp{3.6cm}}
\toprule
对照 & 主要响应 & 解释作用\\
\midrule
解析热核/细FD & 137组有限配对，0.0625\,m细FD最大差0.347\,h & 区分热空间离散与其他状态更新影响\\
单混合物CFL & 0.5降至0.25，存帧水平峰速92.394降至53.256\,m/s & 将时间离散响应纳入物性效应解释\\
D-Claw扩域 & 口岸代表目标响应保持$(160,180]$\,s；1800\,s质量变化约$-4.147\%$ & 同时读取边界响应与密度加权质量账\\
全源面ST网格 & 120/60\,m、同31.184\,Mm$^3$，汇流响应$(240.6,260.1]$ / $(221.0,240.7]$\,s & 保持体积，量化空间离散与足迹变化\\
有界夹带 & 新增5.61843\,Mm$^3$；ASCII扣源质量残差约$-1.17\%$ & 对增长物源建立独立源项账\\
\bottomrule
\end{tabularx}
\end{table}

D-Claw提供固相状态和孔压演化，复合相模型提供相厚度、密度及相间作用，夹带实验提供侵蚀床供给量。这些输出扩大了模型可解释的过程集合。同时，共同栅格输入进入不同框架后形成各自原生初态；例如D-Claw参考初始化的厚度归一$L^1$差约13.90\%，支持由67格变为111格。跨框架差异据此分解为初态、物理过程、边界和离散响应，体现了共同实验设计的作用。

本文采用单混合物120\,m作为主要过程展示，以孔压及复合相路线解释额外状态。复杂度收益通过相应场量、配对变化与源项账呈现；全源面与紧凑源、不同格网以及不同组成属于不同情景，分别保留结果。这样可以把清晰的空间过程展示与量化模型评价放在同一篇论文中。

\subsection{预测信息、区间与评价分母}
时间预测按信息类型分工。物理首次过阈引擎传播初态和参数给定下的状态演化；逆速度与状态空间方法利用发布前连续变形；生存模型通过暴露与失稳组学习动态风险。外部监测评价以13个终末失稳组和1个真实右删失站分组，报告固定发布机会及24、48、72小时上的区间覆盖、宽度和概率分数。不同尺度下的排名变化说明模型选择应与监测频率、发布时刻和预测对象共同确定。

概率输出同时描述时域内启动、时域内不启动和资料/适用域带来的未知质量。机制关系的数值适用域通过固定历史风险集比较，以信息界和条件结果表达其影响。条件95\%时间区间、参数情景范围及外部经验覆盖率分别使用相应定义和分母。对于小样本共形校准，秩计算得到的无界区间作为校准结果保留，模型内区间则用于展示预测分布的形状与尺度。

跨问关联采用同一情景标识，将启动条件、释放状态和目标时延连在一起。7个选定传播运行对应4个机械设计和49条目标记录；重复格网与几何用于配对敏感性。绝对目标时刻按式\eqref{eq:common-clock}保留阶段转换项，空间范围按实际释放足迹及厚度阈值定义。由此形成“机制状态—时间预测—给定释放传播”的可复查链条。

\subsection{监测配置与分级响应的研究建议}
由上述结果可以推导出三类互补的监测和信息配置。第一，区域背景层持续汇集温湿风、辐射和土层水热，按同季节背景及多窗口更新因素族状态；这有利于识别异常过程的持续性。第二，源区状态层优先观测高频变形、裂隙深度温度和水头/排水状态，并校准各设备时钟；它们直接对应机械启动方程及可辨识性分析中的关键状态。第三，沟谷传播层围绕入河位置、汇流段、桥位和设施分别设置响应对象，使传播时钟与工程对象保持一致。

\begin{table}[htbp]
\centering\small
\caption{由模型结构导出的信息配置与响应分工。表中为本文研究建议，阈值配置以具体监测点、工程对象和信息时效为依据。}\label{tab:common-response}
\begin{tabularx}{\linewidth}{p{2.35cm}Y Y}
\toprule
信息层 & 优先观测与更新 & 对应研究输出\\
\midrule
环境背景 & 多高度温湿风、辐射、土层水热、质量轨 & 因素族背景位置、持续时间与排名变化\\
源区状态 & 变形加速、裂隙温度、水压、结构与物源 & 条件稳定裕度、首次过阈、时间分布\\
传播与对象 & 释放状态、河谷响应、桥位和设施几何 & 相对到达括号、流深、影响足迹\\
信息发布 & 校准时钟、起报/有效/可用时刻、实际延迟 & 同发布机会的24/48/72\,h评价\\
\bottomrule
\end{tabularx}
\end{table}

响应组织上，环境背景用于持续跟踪，源区变形与温压用于识别机械状态，整体释放后的沟谷信号用于更新下游响应时窗。各层保留其独立信息来源，再按同一事件和情景关联。本文给出的监测优先级来自模型参数对输出的敏感性及过程可辨识性，可为后续观测设计和防灾策略提供量化依据。
\AIcredit
